\subsection{The exact dual obstruction of the original first five-primary gate}

\begin{result}{E.206}{Proved}
{Poitou--Tate identifies the first-gate cokernel and its adjusted Kummer basis}
Retain a fixed legal minimal-shell pair $q=q_0,r=r_0$, its common
Tate frames, the matrix $M$, and the constant vector $D=(D_q,D_r)$
of E.202.  Put
\[
 S_0=\{2,3,5,q,r,\infty\},\qquad
 V=\mathbb F_5(2),\qquad V^D=\mathbb F_5(3).
\]
Here $V^D=\operatorname{Hom}(V,\mu_5)$ is literally
$\mathbb F_5(-1)$; the displayed identification with
$\mathbb F_5(3)$ uses the trivialization of $\mathbb F_5(4)$
specified by the common Tate generator $\zeta^{\otimes4}$.
In these coordinates evaluation sends
$\zeta^{\otimes2}\otimes\zeta^{\otimes3}$ to $\zeta$.
There are unique classes $w_q,w_r\in H^1(G_{\mathbb Q,S_0},V^D)$
whose tame residue vectors at $q,r$ are $(1,0),(0,1)$ and whose
restriction at five is zero.  Let $\mathcal S^D$ consist of the
classes $\xi$ in this global group satisfying
\[
 \xi|_{G_{\mathbb Q_5}}=0,\qquad
 (\xi(T_q),\xi(F_q))\in\mathbb F_5(24,\alpha),\qquad
 (\xi(T_r),\xi(F_r))\in\mathbb F_5(24,\beta).
 \tag{E.206.1}
\]
The finite local pairings identify
\[
 \operatorname{coker}(M)^*\simeq\mathcal S^D,
 \qquad
 \mathcal S^D=\{xw_q+yw_r:(x,y)^{\mathsf T}\in\ker M^{\mathsf T}\}.
 \tag{E.206.2}
\]
Under this identification the affine obstruction class pairs as
\[
          \langle[D],xw_q+yw_r\rangle=xD_q+yD_r.
 \tag{E.206.3}
\]
Consequently the first system is consistent if and only if this
quantity vanishes for every class in $\mathcal S^D$.
All the local conditions and the possible localization kernel are
accounted for below.  This criterion does not prove its own vanishing.
\end{result}

\begin{proof}
\emph{The fixed-source dual line and its localization at five.}
Let $K=\mathbb Q(\zeta)$ as in E.202 and choose
$\epsilon=(1+\sqrt5)/2$.  The square root is chosen so that
$\sigma_2$ changes its sign.  Kummer descent identifies
\[
 H^1(G_{\mathbb Q,\{5,\infty\}},\mathbb F_5(3))
 \simeq
 \bigl(\mathcal O_K[1/5]^\times/
              \mathcal O_K[1/5]^{\times5}\bigr)^{\chi^{-2}}.
 \tag{E.206.4}
\]
There is no ideal-class contribution, since the class number is one
by E.202.  The ordinary unit rank is one
\citep[Theorem5.1, p.~85]{MilneANT2020}.  The valuation at the
unique prime above five has trivial Galois character, and the
five-primary torsion units have character $\chi$.  Neither
contributes to the character $\chi^{-2}=\chi^2$.
Moreover $\sigma_2\epsilon=-\epsilon^{-1}$, so the Kummer
class of $\epsilon$ has precisely character $\chi^2$.

This class is nonzero even locally at five.  Normalize the valuation
of $K_5=\mathbb Q_5(\zeta)$ by $v(1-\zeta)=1$, so $v(5)=4$.
Then
\[
 \epsilon^4-1=(5+3\sqrt5)/2,\qquad v(\epsilon^4-1)=2.
 \tag{E.206.5}
\]
If a local unit $u$ were a fifth power, say $u=v_0^5$, then
$v_0^4\equiv1$ modulo the maximal ideal, because the residue
field is $\mathbb F_5$.  Every term in the binomial expansion
$(1+t)^5-1$ for $v(t)\ge1$ has valuation at least five.
Thus $v(u^4-1)\ge5$ (or the expression is zero), contradicting
E.206.5 for $u=\epsilon$.
The unit-rank bound now proves that E.206.4 is exactly the line
spanned by the descent of $\kappa_5(\epsilon)\otimes\zeta^2$.

Local duality gives zero $H^0$ and $H^2$ at $\mathbb Q_5$ for
both $V$ and $V^D$: their residual characters $\chi^2,\chi^3$
are nontrivial, and they are each other's Tate duals.
The local Euler characteristic
\citep[I.2.8, p.~31]{CohMilneADT} therefore gives
\[
 \dim H^1(\mathbb Q_5,V)=\dim H^1(\mathbb Q_5,V^D)=1.
\]
The nonzero localization of $\epsilon$ proves that the map
from E.206.4 to $H^1(\mathbb Q_5,V^D)$ is an isomorphism.
This is an integral arithmetic calculation of the relevant
finite line, not an appeal to rational rank alone.

\emph{An explicit basis with its five-adic component removed.}
Use the prime generators $\pi_p$ from E.202, for $p=q,r$, and put
\[
 d_p=\prod_{a=1}^4\sigma_a(\pi_p)^{a^2}.
 \tag{E.206.6}
\]
Its Kummer class has character $\chi^{-2}=\chi^2$, by the
same exponent reindexing as E.202.10.  It has valuation one at
the chosen prime over $p$ and no other rational prime support
apart from $p$.  There is a unique $\lambda_p\in\mathbb F_5$
such that its local Kummer class at five is
$\lambda_p[\epsilon]$ in the $\chi^2$ eigenspace.  Define
\[
                    e_p=d_p\epsilon^{-\lambda_p}.
 \tag{E.206.7}
\]
The unique descent of $\kappa_5(e_p)\otimes\zeta^2$ is $w_p$.
It has the asserted residue vector and zero localization at five.
Changing a prime generator changes its weighted Kummer class
only by the global $\chi^2$ unit line, which E.206.7 removes.
Thus the resulting cohomology classes are independent of that
choice.  The same conclusion follows from their uniqueness:
any difference has zero residues, hence belongs to E.206.4,
where localization at five is injective.

More explicitly, Kummer theory, class number one and the
$S_0$-unit valuation sequence show that
\[
 H^1(G_{\mathbb Q,S_0},V^D)
   =\mathbb F_5[\epsilon]\oplus\mathbb F_5w_q\oplus\mathbb F_5w_r.
 \tag{E.206.8}
\]
The primes two and three add no $\chi^2$ valuation eigenspace:
each has Frobenius square $-1$ modulo five.  The primes $q,r$
split completely in $K$ and each add exactly one such eigenspace.
In particular the kernel of localization at five is exactly
$\mathbb F_5w_q\oplus\mathbb F_5w_r$.

\emph{The local spaces, including all real terms.}
At $p=2,3$ both $H^1(\mathbb Q_p,V)$ and
$H^1(\mathbb Q_p,V^D)$ are zero.  Indeed the unramified term
has scalar $p^s-1$ and the tame term has scalar $p^{s-1}-1$
for weight $s$; both are units modulo five for $s=2,3$ and
$p=2,3$.  At $q,r$ the two characters are trivial modulo five,
and evaluation at $(T_p,F_p)$ identifies each local $H^1$
with $\mathbb F_5^2$.  The real decomposition group has order
two, invertible on these coefficient groups.  All of its Tate
cohomology groups, including the degree-zero endpoint used
in global duality, vanish by averaging.  There is therefore
no omitted real correction in this calculation.

Let
\[
 \mathcal R=\bigoplus_{v\in S_0}H^1(\mathbb Q_v,V),\qquad
 \mathcal R^D=\bigoplus_{v\in S_0}H^1(\mathbb Q_v,V^D).
\]
Both spaces have dimension five, contributed by $5,q,r$ as
$1+2+2$.  The global localization maps are injective.
For $V$ this follows from the residue isomorphism E.202.3.
For $V^D$, a class locally zero at $q,r$ has zero residues,
so lies in E.206.4; local vanishing at five then kills it.
Thus no quotient by an unaccounted-for localization kernel
is required in E.206.2.

\emph{The exact finite Poitou--Tate map.}
For a finite module with order invertible outside $S_0$,
\citet[I.4.9--4.10(b), pp.~56--57]{CohMilneADT} gives exactness
of the segment
\[
 H^1(G_{\mathbb Q,S_0},V)\longrightarrow\mathcal R
 \longrightarrow H^1(G_{\mathbb Q,S_0},V^D)^*,
 \tag{E.206.9}
\]
where the second map is the sum of the local cup-product
invariant pairings.  Its dual statement says that the two
global localization images are exact annihilators.
The archimedean convention in that source uses Tate cohomology,
whose vanishing was checked explicitly above.  All spaces here
are finite; no inverse-limit assertion is being substituted
for E.206.9.  As an independent dimension check, the images
have dimensions two and three by E.202.3 and E.206.8, and
local duality pairs two spaces of dimension five.

Fix the standard arithmetic-Frobenius normalization of the
local invariant.  With the common Tate coordinate, its value
at $p=q,r$ is
\[
 \langle z,\xi\rangle_p
       =z(F_p)\xi(T_p)-z(T_p)\xi(F_p)\in\mathbb F_5.
 \tag{E.206.10}
\]
To check the sign, the cup of the unramified character taking
$F_p$ to one with the Kummer class of the rational uniformizer
has invariant $1/5$; multiplication by five gives one.
Indeed the local Artin map sends that uniformizer to arithmetic
Frobenius, and the invariant of this cup is the value of the
character on its Artin image
\citep[III, Proposition3.6 and Lemma3.7, pp.~109--110;
\S4, Steps2--3 and Proposition4.3, pp.~112--113]{RootMilneCFT}.
The reversed cup has the negative value.  Their self-cups
vanish at odd order, giving E.206.10 by bilinearity.
Changing every local invariant to its opposite would change
both sides of the global pairing by the same harmless sign;
we keep E.206.10 throughout.

Impose the following homogeneous local conditions on $V$:
at five allow the entire local group, while at $q,r$ impose
\[
 z(F_q)-\alpha z(T_q)/24=0,\qquad
 z(F_r)-\beta z(T_r)/24=0.
 \tag{E.206.11}
\]
The quotient of $\mathcal R$ by these conditions is
$\mathbb F_5^2$, and global localization to it is precisely
$M$ in the residue basis $z_q,z_r$.
The annihilator conditions on the dual are zero at five and
the two lines E.206.1, by E.206.10.  Exact annihilation in
E.206.9, together with the just-proved injectivity, gives the
first isomorphism E.206.2.  This proves the exact duality
interpretation of this particular system, with its source
and every local condition specified.

\emph{The transpose identities and the affine pairing.}
Global reciprocity applied to $z_i\cup w_j$ has no term at
five because $w_j$ is locally zero there, none at two or
three by the local vanishing, and none at infinity.  Outside
$S_0$ both classes are unramified and their cup is zero,
since the finite residue field has cohomological dimension
one.  Applying E.206.10 at $q,r$ consequently gives
\[
 \begin{aligned}
 w_q(F_q)&=z_q(F_q),&\quad w_r(F_q)&=z_q(F_r),\\
 w_q(F_r)&=z_r(F_q),&\quad w_r(F_r)&=z_r(F_r).
 \end{aligned}
 \tag{E.206.12}
\]
For $\xi=xw_q+yw_r$ the residues are $x,y$.
Its two line conditions E.206.1, after substitution of
E.206.12, are exactly $M^{\mathsf T}(x,y)^{\mathsf T}=0$.
This proves the second equality in E.206.2.
Represent the target vector $D$ by local $V$ classes with
zero tame values and Frobenius values $D_q,D_r$ and with
zero component at five.  Equation E.206.10 then gives
exactly $xD_q+yD_r$.  Replacing the representative by a
homogeneous local condition or a global correction does
not change this number for $\xi\in\mathcal S^D$.
This proves E.206.3 and the consistency criterion.

The Kummer expressions make E.206.12 an explicit quintic
reciprocity identity: in the uniformizer-fixing frames of
E.202, replace $w_p(F_p)$ by $\ell_p(e_p/p)$ and its
off-diagonal values by $\ell_q(e_r),\ell_r(e_q)$.
The result is adjointness of the two different Tate-weight
matrices.  It does not identify those two weights globally,
force a diagonal entry to vanish, or evaluate the affine
pairing in E.206.3.  In particular the weight-one Kummer
classes $a,b$ only specify the two displayed local lines;
there is no claimed global weight-three linear combination
of $a,b$ that would annihilate that pairing automatically.
\end{proof}
